% GENERATED FILE. Edit canonical lessons, then run npm run generate:books.
% canonical-source-hash: fnv1a64:e887fdb76f1f5755
% canonical-lessons: RU-C217-tupoy, RU-C217-gladkiy, RU-C217-prokhladnyy, RU-C217-svezhiy, RU-C217-sovremennyy

\chapter{Blunt, Smooth, Cool, Fresh, Modern}
\label{ch:ru-blunt-smooth-cool-fresh-modern}
\hlchaptermodality{217}

\begin{chapteropening}
\textbf{By the end of this chapter:} \emph{I can say blunt, smooth, cool, fresh and modern.}

Complete the last lesson of chapter 217: I can say blunt, smooth, cool, fresh and modern.
\end{chapteropening}

\section[\texorpdfstring{\ru{тупой} (tupóy)}{тупой (tupóy)}]{\ru{тупой} (tupóy) — blunt}

\label{lesson:RU-C217-tupoy}



\begin{quote}
Before the new one: say the Russian for sad, then the Russian for strict.
\end{quote}

\subsection*{You'll want to know: \ru{тупой}}
\textbf{\ru{тупой}} (\emph{tupóy}) — \textquotedblleft{}blunt\textquotedblright{}.

\begin{grammarlens}[title={What you've built}]
One more word for this chapter.
\end{grammarlens}

\subsection*{Your turn}
\begin{itemize}
\raggedright
  \item \emph{Say it:} \emph{tupóy}
  \item \emph{Say it:} \emph{tupóy}, once more
  \item \emph{Say it:} say \emph{tupóy}
  \item \emph{Recall:} say the Russian for sad, then the Russian for strict, then say \emph{tupóy} again
\end{itemize}

\begin{tcolorbox}[breakable,colback=teal!4,colframe=teal!35!black,title={Before you move on}]
What does \ru{тупой} mean? (\textquotedblleft{}Blunt\textquotedblright{}.) Say it once more.
\end{tcolorbox}
\section[\texorpdfstring{\ru{гладкий} (gládkiy)}{гладкий (gládkiy)}]{\ru{гладкий} (gládkiy) — smooth}

\label{lesson:RU-C217-gladkiy}



\begin{quote}
Before the new one: say the Russian for strict, then the Russian for blunt.
\end{quote}

\subsection*{You'll want to know: \ru{гладкий}}
\textbf{\ru{гладкий}} (\emph{gládkiy}) — \textquotedblleft{}smooth\textquotedblright{}.

\begin{grammarlens}[title={What you've built}]
One more word for this chapter.
\end{grammarlens}

\subsection*{Your turn}
\begin{itemize}
\raggedright
  \item \emph{Say it:} \emph{gládkiy}
  \item \emph{Say it:} \emph{gládkiy}, once more
  \item \emph{Say it:} say \emph{gládkiy}
  \item \emph{Recall:} say the Russian for strict, then the Russian for blunt, then say \emph{gládkiy} again
\end{itemize}

\begin{tcolorbox}[breakable,colback=teal!4,colframe=teal!35!black,title={Before you move on}]
What does \ru{гладкий} mean? (\textquotedblleft{}Smooth\textquotedblright{}.) Say it once more.
\end{tcolorbox}
\section[\texorpdfstring{\ru{прохладный} (prokhládnyy)}{прохладный (prokhládnyy)}]{\ru{прохладный} (prokhládnyy) — cool}

\label{lesson:RU-C217-prokhladnyy}



\begin{quote}
Before the new one: say the Russian for blunt, then the Russian for smooth.
\end{quote}

\subsection*{You'll want to know: \ru{прохладный}}
\textbf{\ru{прохладный}} (\emph{prokhládnyy}) — \textquotedblleft{}cool\textquotedblright{}.

\begin{grammarlens}[title={What you've built}]
One more word for this chapter.
\end{grammarlens}

\subsection*{Your turn}
\begin{itemize}
\raggedright
  \item \emph{Say it:} \emph{prokhládnyy}
  \item \emph{Say it:} \emph{prokhládnyy}, once more
  \item \emph{Say it:} say \emph{prokhládnyy}
  \item \emph{Recall:} say the Russian for blunt, then the Russian for smooth, then say \emph{prokhládnyy} again
\end{itemize}

\begin{tcolorbox}[breakable,colback=teal!4,colframe=teal!35!black,title={Before you move on}]
What does \ru{прохладный} mean? (\textquotedblleft{}Cool\textquotedblright{}.) Say it once more.
\end{tcolorbox}
\section[\texorpdfstring{\ru{свежий} (svézhiy)}{свежий (svézhiy)}]{\ru{свежий} (svézhiy) — fresh}

\label{lesson:RU-C217-svezhiy}



\begin{quote}
Before the new one: say the Russian for smooth, then the Russian for cool.
\end{quote}

\subsection*{You'll want to know: \ru{свежий}}
\textbf{\ru{свежий}} (\emph{svézhiy}) — \textquotedblleft{}fresh\textquotedblright{}.

\begin{grammarlens}[title={What you've built}]
One more word for this chapter.
\end{grammarlens}

\subsection*{Your turn}
\begin{itemize}
\raggedright
  \item \emph{Say it:} \emph{svézhiy}
  \item \emph{Say it:} \emph{svézhiy}, once more
  \item \emph{Say it:} say \emph{svézhiy}
  \item \emph{Recall:} say the Russian for smooth, then the Russian for cool, then say \emph{svézhiy} again
\end{itemize}

\begin{tcolorbox}[breakable,colback=teal!4,colframe=teal!35!black,title={Before you move on}]
What does \ru{свежий} mean? (\textquotedblleft{}Fresh\textquotedblright{}.) Say it once more.
\end{tcolorbox}
\section[\texorpdfstring{\ru{современный} (sovreménnyy)}{современный (sovreménnyy)}]{\ru{современный} (sovreménnyy) — modern}

\label{lesson:RU-C217-sovremennyy}



\begin{quote}
Before the new one: say the Russian for cool, then the Russian for fresh.
\end{quote}

\subsection*{You'll want to know: \ru{современный}}
\textbf{\ru{современный}} (\emph{sovreménnyy}) — \textquotedblleft{}modern\textquotedblright{}.

\begin{grammarlens}[title={What you've built}]
One more word for this chapter.
\end{grammarlens}

\subsection*{Your turn}
\begin{itemize}
\raggedright
  \item \emph{Say it:} \emph{sovreménnyy}
  \item \emph{Say it:} \emph{sovreménnyy}, once more
  \item \emph{Say it:} say \emph{sovreménnyy}
  \item \emph{Recall:} say the Russian for cool, then the Russian for fresh, then say \emph{sovreménnyy} again
\end{itemize}

\begin{tcolorbox}[breakable,colback=teal!4,colframe=teal!35!black,title={Before you move on}]
What does \ru{современный} mean? (\textquotedblleft{}Modern\textquotedblright{}.) Say it once more.
\end{tcolorbox}
